\documentclass[11pt]{article}
\usepackage[margin=1in]{geometry}
\usepackage{amsmath,amssymb,amsthm}
\usepackage{mathtools}
\usepackage{hyperref}

\newtheorem{theorem}{Theorem}
\newtheorem{proposition}[theorem]{Proposition}
\newtheorem{lemma}[theorem]{Lemma}
\newtheorem{corollary}[theorem]{Corollary}
\theoremstyle{definition}
\newtheorem{definition}[theorem]{Definition}

\newcommand{\R}{\mathbb{R}}
\newcommand{\C}{\mathbb{C}}
\newcommand{\N}{\mathbb{N}}
\newcommand{\Cpls}{\C^+}
\newcommand{\Sprime}{\mathcal{S}'}
\DeclareMathOperator{\Real}{Re}
\DeclareMathOperator{\Imag}{Im}
\DeclareMathOperator{\supp}{supp}

\title{The Outer-Phase Identity for the Warped Pullback Spectrum}
\author{Stephen Crowley}
\date{}

\begin{document}
\maketitle

\begin{abstract}
Let $H(x) := Z(\theta^{-1}(x))/\sqrt{\theta'(\theta^{-1}(x))}$ be the warped pullback of the Hardy $Z$-function, with $\theta$ the Riemann--Siegel phase. By Theorem~B of the companion document (``Three Missing Theorems''), the windowed modulus $|\widehat H_T(\xi)|^2/(\theta(T) - X_0)$ converges to a nonnegative, compactly-supported spectral density $S(\xi)$ with explicit mode-sum form on $(-1,1)$, and the whitened differential factorizes as $d\widehat H = \sqrt{S(\xi)}\,e^{i\Phi(\xi)}\,d\xi$ in $\Sprime([-1+\delta, 1-\delta])$ with $\Phi(\xi) = \arg\sum_{n,\sigma}|\mathcal A_{n,\sigma}(\xi)|e^{i\phi_{n,\sigma}(\xi)}$. Define the analytic signal $H_+(z) := \int_0^1 e^{i\xi z}\,d\widehat H(\xi)$. We reduce the question ``$H_+$ is outer in $\Cpls$'' to a single explicit identity on $(-1+\delta, 1-\delta)$:
\[
\Phi(\xi) \equiv \mathcal H\!\left[\tfrac12 \log S\right](\xi) \quad \pmod{\pi},
\]
where $\mathcal H$ is the Hilbert transform on $(-1,1)$. This is the Krein--Szeg\H{o} outer-phase condition specialized to the explicit phase produced by Theorem~B.
\end{abstract}

\section{Inputs}

\begin{definition}[Warped pullback]
$H(x) := Z(\theta^{-1}(x))/\sqrt{\theta'(\theta^{-1}(x))}$, $x \geq X_0 := \theta(T_0)$, $T_0 \geq 200$.
\end{definition}

\begin{definition}[Windowed Fourier transform]
$\widehat H_T(\xi) := \int_{X_0}^{\theta(T)} H(x)\, e^{-i\xi x}\, dx$.
\end{definition}

\paragraph{Inputs from prior results (cited from companion documents).}

\textbf{(Band-limit, \texttt{BandLimitedness.tex}, Theorem~9).} $|\widehat H_T(\xi)|^2/(\theta(T) - X_0) \to 0$ for every fixed $|\xi| > 1$, locally uniformly.

\textbf{(Theorem~B of \texttt{three\_missing\_theorems}).} For almost every $\xi \in (-1, 1)$,
\begin{equation}\label{eq:B1}
\lim_{T \to \infty}\frac{|\widehat H_T(\xi)|^2}{\theta(T) - X_0} = S(\xi) := \frac{1}{(2\pi)^2}\sum_{n \geq 1,\, \sigma \in \{\pm\}}|\mathcal A_{n,\sigma}(\xi)|^2,
\end{equation}
with $\mathcal A_{n,\sigma}(\xi) = n^{-1/2}\sqrt{2\pi/|\theta''(t^*_{n,\sigma,\xi})|}\,e^{i\pi/4}$ the stationary-phase amplitude and $t^*_{n,\sigma,\xi} := (\theta')^{-1}(\sigma\log n/(\sigma - \xi))$. Moreover, on $\Sprime([-1+\delta, 1-\delta])$,
\begin{equation}\label{eq:B4}
d\widehat H(\xi) = \sqrt{S(\xi)}\,e^{i\Phi(\xi)}\,d\xi, \qquad \Phi(\xi) := \arg\!\sum_{n,\sigma}|\mathcal A_{n,\sigma}(\xi)|\,e^{i\phi_{n,\sigma}(\xi)},
\end{equation}
where $\phi_{n,\sigma}(\xi)$ is the stationary-phase phase of mode $(n, \sigma)$.

\section{The analytic signal and its $H^\infty(\Cpls)$ factorization}

\begin{definition}[Analytic signal]
\[
H_+(z) := 2\int_{[0,1]} e^{i\xi z}\, d\widehat H(\xi) = 2\int_0^1 e^{i\xi z}\sqrt{S(\xi)}\,e^{i\Phi(\xi)}\,d\xi, \qquad z \in \C.
\]
\end{definition}

\begin{lemma}[Basic properties]\label{lem:basic}
$H_+$ is entire; $|H_+(z)| \leq 2\int_0^1 \sqrt{S(\xi)}\, d\xi < \infty$ on $\overline{\Cpls}$; $|H_+(iy)| \to 0$ as $y \to +\infty$; on $\R$, $H_+(u) = H(u) + i\mathcal H[H](u)$.
\end{lemma}

\begin{proof}
Absolute convergence of the integral for all $z \in \C$ from compact support of the integrand. Bound in $\Cpls$ by $|e^{i\xi z}| = e^{-\xi \Imag z} \leq 1$ for $\xi \geq 0$, $\Imag z \geq 0$. Decay on $i\R^+$ by dominated convergence: $e^{-\xi y} \to 0$ on $(0, 1]$. Reality identity from $\widehat{H_+} = 2\mathbf 1_{[0,1]}\widehat H$ combined with $\widehat{H}(-\xi) = \overline{\widehat H(\xi)}$.
\end{proof}

\begin{lemma}[$H^\infty(\Cpls)$ canonical factorization]\label{lem:canonical}
$H_+ \in H^\infty(\Cpls)$ admits the canonical factorization $H_+ = B \cdot S_\infty \cdot F$ in $\Cpls$, with $B$ a Blaschke product over zeros of $H_+$ in $\Cpls$, $S_\infty$ a singular inner factor supported at infinity, and $F$ outer. Since $H_+$ is entire with boundary values controlled by Lemma~\ref{lem:basic}, $S_\infty \equiv 1$ (no singular inner at infinity: Lemma~\ref{lem:basic} gives sub-exponential decay of $\log|H_+(iy)|$ relative to $y$). Hence
\[
H_+(z) = B(z)\, F(z), \qquad z \in \Cpls.
\]
\end{lemma}

\begin{proof}
Standard canonical factorization for $H^\infty$ of the upper half-plane (Garnett, \emph{Bounded Analytic Functions}, Ch.~II, Thm.~4.6). Triviality of $S_\infty$: by Lemma~\ref{lem:basic}, $|H_+(iy)| \to 0$ but without exponential rate; indeed $\log|H_+(iy)|/y \to 0$ as $y \to \infty$, which is the criterion for trivial singular inner factor at infinity for half-plane $H^\infty$.
\end{proof}

\begin{corollary}\label{cor:outer-iff-BtrivialB}
$H_+ \in$ Hermite--Biehler class on $\Cpls$ (i.e., $H_+$ has no zeros in $\Cpls$) if and only if $B \equiv 1$, i.e., $H_+ = F$ is outer in $\Cpls$.
\end{corollary}

\section{Krein--Szeg\H{o} outer function from the spectral density}

\begin{definition}[Szeg\H{o} outer function of $S$ on the band]
Let $\delta \in (0,1)$ small. For $\xi \in (-1+\delta, 1-\delta)$, assume $S(\xi) > 0$ (positivity of the spectral density on the interior, from Theorem~B mode-sum structure). Define
\[
F_S(z) := \exp\!\left(\frac{1}{2\pi i}\int_{-1+\delta}^{1-\delta}\frac{1 + \xi z}{\xi - z}\cdot\frac{\log S(\xi)}{1 + \xi^2}\,d\xi\right), \qquad z \in \Cpls.
\]
This is the Krein--Szeg\H{o} outer function with $|F_S(\xi)|^2 = S(\xi)$ a.e.\ on $(-1+\delta, 1-\delta)$ and $F_S \in H^2(\Cpls)$ outer.
\end{definition}

\begin{lemma}[Boundary phase of $F_S$]\label{lem:szego-phase}
The boundary values of $F_S$ on $(-1+\delta, 1-\delta)$ satisfy
\[
F_S(\xi + i0^+) = \sqrt{S(\xi)}\cdot e^{i\Psi(\xi)}, \qquad \Psi(\xi) := \mathcal H_\delta\!\left[\tfrac12 \log S\right]\!(\xi),
\]
where $\mathcal H_\delta$ is the Hilbert transform on $(-1+\delta, 1-\delta)$ with kernel $(1 + \xi \eta)/((\xi - \eta)(1 + \eta^2))$ (the finite-interval Plemelj kernel corresponding to the Schwarz integral).
\end{lemma}

\begin{proof}
Apply the Plemelj formula to the Schwarz integral defining $\log F_S$. The Schwarz kernel $(1 + \xi z)/(\xi - z)\cdot 1/(1+\xi^2)$ decomposes on the real axis boundary as $\pi i\,\delta(\xi - \eta) + \text{p.v.}\,(1+\xi\eta)/((\xi-\eta)(1+\eta^2))$. Taking $\Real$ and $\Imag$ of $\log F_S(\eta + i0^+)$ gives $\Real \log F_S = \tfrac12 \log S$ and $\Imag \log F_S = \mathcal H_\delta[\tfrac12 \log S]$.
\end{proof}

\section{The outer-phase identity}

\begin{theorem}[Outer-phase identity]\label{thm:main}
$H_+$ is outer in $\Cpls$ if and only if
\begin{equation}\label{eq:identity}
\Phi(\xi) \equiv \mathcal H_\delta\!\left[\tfrac12\log S\right]\!(\xi) \pmod{\pi}, \qquad \text{a.e. } \xi \in (-1+\delta, 1-\delta),
\end{equation}
for every $\delta \in (0,1)$, where $\Phi$ is the whitened phase of \eqref{eq:B4} and $\mathcal H_\delta$ is the Hilbert transform of Lemma~\ref{lem:szego-phase}.
\end{theorem}

\begin{proof}
Two boundary-value functions of $H^2(\Cpls)$ functions agree as boundary values (up to a sign in the boundary argument) iff the functions agree up to sign. By Lemma~\ref{lem:basic} and \eqref{eq:B4}, $H_+$ has boundary values on $\R$ whose restriction to $(-1+\delta, 1-\delta)$ in the $d\widehat H$-sense is $\sqrt{S(\xi)}e^{i\Phi(\xi)}\,d\xi$. The Krein--Szeg\H{o} outer function $F_S$ has boundary values $\sqrt{S(\xi)}e^{i\Psi(\xi)}\,d\xi$ (Lemma~\ref{lem:szego-phase}) with $\Psi = \mathcal H_\delta[\tfrac12\log S]$.

\emph{Forward.} If $H_+$ is outer, then $B \equiv 1$ in Lemma~\ref{lem:canonical}, so $H_+ = F$ coincides with the outer function having $|F(\xi)|^2 = S(\xi)$ a.e., which is $F_S$ up to a unimodular constant. Matching boundary phases gives \eqref{eq:identity}.

\emph{Converse.} Suppose \eqref{eq:identity} holds. Then on $(-1+\delta, 1-\delta)$, $d\widehat H(\xi) = \sqrt{S(\xi)}e^{i\mathcal H_\delta[\tfrac12\log S](\xi)}\,d\xi$ equals the boundary value of $F_S\cdot d\xi$. Since $\delta > 0$ is arbitrary, the identity extends to $(-1, 1)$ a.e. Hence
\[
H_+(z) = 2\int_0^1 e^{i\xi z}\sqrt{S(\xi)}e^{i\mathcal H[\tfrac12\log S](\xi)}\,d\xi = 2 F_S(z)\quad \text{(up to sign / unimodular constant)},
\]
which is outer in $\Cpls$ by construction.
\end{proof}

\begin{corollary}\label{cor:RH-equivalence}
$H$ has only real zeros (Riemann Hypothesis in warped coordinates) if and only if the identity \eqref{eq:identity} holds.
\end{corollary}

\begin{proof}
$H$ has only real zeros $\iff$ $H_+$ is outer in $\Cpls$ (Hermite--Biehler: Corollary~\ref{cor:outer-iff-BtrivialB} combined with the symmetrization $\tilde H = \tfrac12(H_+ + H_+^*)$ having zeros exactly at real zeros of $H$) $\iff$ \eqref{eq:identity} (Theorem~\ref{thm:main}).
\end{proof}

\section{What \eqref{eq:identity} says concretely}

Both sides of \eqref{eq:identity} are explicit functions on $(-1+\delta, 1-\delta)$:

\textbf{LHS.} $\Phi(\xi) = \arg\sum_{n \geq 1,\sigma \in \{\pm\}}|\mathcal A_{n,\sigma}(\xi)|\,e^{i\phi_{n,\sigma}(\xi)}$. Both $|\mathcal A_{n,\sigma}(\xi)|$ and $\phi_{n,\sigma}(\xi)$ are explicit functions of $\xi$ through the stationary-phase locus $t^*_{n,\sigma,\xi}$ and the Riemann--Siegel phase $\theta$. For mode $(n, +)$: $t^*_{n,+,\xi} = (\theta')^{-1}(\log n/(1 - \xi))$, $|\mathcal A_{n,+}(\xi)|^2 = 2\pi\log n / (n(1-\xi)^2 \theta''(t^*_{n,+,\xi}))$, and $\phi_{n,+}(\xi) = (1-\xi)\theta(t^*_{n,+,\xi}) - t^*_{n,+,\xi}\log n - \pi/4$.

\textbf{RHS.} $\mathcal H_\delta[\tfrac12\log S]$ with $S$ the same mode-sum of squared amplitudes divided by $(2\pi)^2$.

The identity \eqref{eq:identity} says: \emph{the argument of the mode-sum $\sum|\mathcal A_{n,\sigma}|e^{i\phi_{n,\sigma}}$ equals the Hilbert transform of one-half the logarithm of the sum of squared amplitudes.} This is a nontrivial identity between two computable functions.

\section{Status}

Theorem~\ref{thm:main} is a clean reformulation: the ``only real zeros'' question for $H$ becomes the single-identity question \eqref{eq:identity}. Neither direction is proved in this document. The LHS and RHS are both fully explicit from Theorem~B and the Krein--Szeg\H{o} apparatus. What remains is:

\begin{enumerate}
\item Verification of \eqref{eq:identity} on numerical intervals (computable to any precision from the mode-sum formulas; provides evidence but not proof).
\item Analytic proof of \eqref{eq:identity} as an identity of Fourier-analytic objects, exploiting the Riemann--Siegel structure.
\end{enumerate}

The identity \eqref{eq:identity} is the target. It is what must be proven to close the outer-in-$\Cpls$ route.

\begin{thebibliography}{9}
\bibitem{Garnett} J.~B.~Garnett, \emph{Bounded Analytic Functions}, Revised 1st ed., Graduate Texts in Mathematics 236, Springer, 2007.
\bibitem{Koosis} P.~Koosis, \emph{Introduction to $H_p$ Spaces}, 2nd ed., Cambridge, 1998.
\bibitem{Levin} B.~Ya.~Levin, \emph{Distribution of Zeros of Entire Functions}, AMS Translations of Mathematical Monographs 5, 1980.
\bibitem{ThreeMissing} ``Three Missing Theorems,'' companion document, 2026 (Theorem~B: diagonal-dominance identity for the whitened differential).
\bibitem{BandLimit} ``Band-Limitedness of the Warped Pullback of the Hardy $Z$-Function,'' companion document, 2026.
\end{thebibliography}

\end{document}
